%%%%%%%%%%%%%%%%%%%%%%%%%%%%%%%%%%%%%%%%%%%%%%%%%%%%%%%%%%%%%%%%%%%%%%%%%%%%%%%
\section{Limitations}
\label{sec:limits}

\emph{The simulated activity.} The quasi-periodic FF$'$ model is a common description of
spot-driven variability, but it is an assumption, and so is every prior on its
parameters. Most consequential is how activity imprints on the indicators: we vary the
amplitude, the lag and the share of indicator-free cases (Sect.~\ref{sec:sens}), but the
fifth of cases without any indicator signal is an arbitrary choice, and the classifiers
learn the coupling we assumed. The kernel also puts little power in high harmonics of
the rotation period.

\emph{The noise model.} Permuting residuals within eras keeps the error distribution,
its tails and the joint RV--indicator noise, but removes correlated noise on timescales
longer than a night, including magnetic cycles and slow instrumental drifts. The
false-alarm rates of Sect.~\ref{sec:fap} are for noise without such correlations.

\emph{The priors.} The fraction of significant peaks that are impostors, and the
classifiers' probabilities, depend on the class mix and the amplitude priors. They
describe these simulations, not the archive.

\emph{The benchmark.} It has \nmBenchBinary\ signals with a verdict and only
\nmBRuleRefN\ significant refuted ones, and its results depend on the window allowed
around a claimed period (Table~\ref{tab:window}). The rotation proxy behaves differently
on real indicators than on simulated ones (Sect.~\ref{sec:bench}). Its labels come from later literature that often
used similar tests, and some of its upheld signals are easy, \nmBEasyN\ with
$K/\sigma_K > 20$. It measures agreement with the literature on the full archive, not the
performance a discoverer would have had.

\emph{Changes after the first pass.} The rule's thresholds were fixed before the
benchmark was run, but the pipeline changed four times afterwards
(Sect.~\ref{sec:changes}), and some thresholds came from the analysis of HD\,297396\,b.

\emph{One instrument.} Every test here uses HARPS data alone. A signal specific to one
star in one spectrograph passes every rung; only a second instrument addresses it.

%%%%%%%%%%%%%%%%%%%%%%%%%%%%%%%%%%%%%%%%%%%%%%%%%%%%%%%%%%%%%%%%%%%%%%%%%%%%%%%
\section{Conclusions}
\label{sec:conc}

\begin{enumerate}
\item Neither common treatment of the jitter gives nominal false-alarm rates on real
      archival noise. Held at its planet-free value, it makes the Baluev approximation
      conservative by a factor of \nmFapFixFactorLo--\nmFapFixFactorHi; rescaled under the
      signal model, liberal by a factor of \nmFapLiberalPerm\ at 1\% and
      \nmFapLiberalPermTenth\ at 0.1\%. On these stars a Baluev FAP of \nmFapEmpOne\%
      corresponds to a real false-alarm rate of 1\%.
\item Evolving activity produces many significant peaks inside the broad lines of its
      power spectrum rather than on exact rotation harmonics; the balance shifts from
      lines to harmonics as spots live longer (\nmLineShort\ against \nmHarmShort\ for
      lifetimes under three rotations, \nmLineLong\ against \nmHarmLong\ above ten). Tests of
      harmonic distance see only the harmonics, and only when the indicators reveal the
      rotation period.
\item No single test is decisive: the best reaches an AUC of \nmAucInd\ and the detection
      statistics \nmAucFap--\nmAucDchi. Combined, the rungs reach \nmAucLog--\nmAucTree, and
      \nmAucNoInd\ without the indicators. How strongly activity shows in the indicators
      matters more than any other property of the simulated activity.
\item On published signals, \nmRefNotDet\ of \nmBenchR\ refuted signals are not
      significant in a blind search of the full archive, though \nmRefDetSingle\ keep
      single-frequency power at their period. The rule with thresholds fixed in advance rejects
      \nmRefFFInd\ of the \nmBRuleRefN\ significant refuted signals on indicators,
      \nmRefFFPeriod\ on period and \nmRefFFPhase\ on phase, accepts \nmBRuleRefAcc, and rejects \nmBRuleUpRej\ of \nmBRuleUpN\
      significant upheld signals; the classifiers rank the significant signals with AUCs
      of \nmBAucTree\ (trees) and \nmBAucLog\ (logistic). These counts shift with the window
      allowed around the claimed period, and the numbers are too small to bound either
      error rate tightly.
\item Run on the archive as it stood when it closed, the ladder's list held \nmTsHits\ signals
      later published as planets, gathered at the top (\nmTsTreeTen\ of the ten ranked
      highest by the trees, against \nmTsExpTen\ expected), but the classifiers did not rank
      them better than the false-alarm probability alone, and \nmTsToiTopTen\ of those
      \nmTsTreeTen\ orbit stars with TESS candidates that observers were already following. Of
      two open signals tested with newer spectra under a rule written in advance, one was
      refuted and the other is undecided under the registered model (about
      $\nmTaHdSigma\sigma$ with a zero-point split chosen after the velocities were displayed).
\item In the AI-assisted analysis that motivated this work, \nmDefN\ defects were found
      in three revisions. Overclaims were the largest category (\nmDefCatOverclaim), and
      transcriptions of the code's own output were few (\nmDefOriCode); an independent
      reviewer that saw the sources but not the history found \nmDefRouteRev\ of them.
\end{enumerate}

\section*{Acknowledgements}

This work is based on data from the HARPS-RVBank \citep{Trifonov2020,Perdelwitz2024},
built from observations collected at the European Southern Observatory and
retrieved from the ESO Science Archive Facility. It made use of NASA's Astrophysics
Data System, the NASA Exoplanet Archive, SIMBAD and VizieR, and of \texttt{numpy}
\citep{Harris2020}, \texttt{scipy} \citep{Virtanen2020}, \texttt{scikit-learn}
\citep{Pedregosa2011} and \texttt{matplotlib} \citep{Hunter2007}. Section~\ref{sec:timesplit}
uses public HARPS, ESPRESSO and NIRPS spectra from the ESO Science Archive Facility,
from ESO programmes 106.21R4.001, 108.222V.001, 110.248C.001, 112.25YG.002,
112.25YG.005, 112.25YG.006, 112.25YG.007, 112.25YG.008, 112.25YG.009, 112.25YG.010 and
115.285G.001.

\emph{Use of generative AI.} The author designed the study and made the analysis
decisions. The code, the literature search for the benchmark, the extraction and coding
of the defect log, and the draft text were produced with the assistance of Claude
(Anthropic), a large language model. They were validated as follows. The package has 40
unit tests, 13 of which compare it with an independent calculation
(Sect.~\ref{sec:verify}). Every simulated case is reproducible from its seed: \nmConsistN\ randomly chosen
cases regenerated with the released code gave features \nmConsistDiff\ to the stored
ones. Every
benchmark label, period and reference, and every entry of the bibliography, was checked
against the publisher's record. A separate instance of the model, given the manuscript,
the code and the outputs but not the conversation that produced them, reviewed the first
draft as a referee; its report and our response are released with the code. The defect
log was extracted and coded by separate instances following a codebook fixed in
advance, with agreement measured between two coders. Every computed result in the
text is generated by the analysis scripts from their outputs; design constants, literature
values and counts typed by hand are listed in the released provenance file. The author takes responsibility for the content.

\section*{Data and code availability}

HARPS-RVBank is available at the CDS (J/A+A/683/A125). The package \code{rvvet}, the
simulation outputs, the false-alarm calibration, the benchmark table with its
references, the frozen rule and the frozen-pipeline results, the defect log with its
codebook, sources and both codings, the referee report, the forward test of
Sect.~\ref{sec:timesplit} (the archive run, the catalogue cross-matches, the discovery ledger
and the prospective test with its velocities), the pinned software environment and every
script used for this paper are released at \url{https://github.com/samsonmaximus/rvvet} and archived at Zenodo
(\href{https://doi.org/10.5281/zenodo.23249395}{doi:10.5281/zenodo.23249395}).
The benchmark of published signals and the January 2022 list of Sect.~\ref{sec:timesplit}
are also ancillary files of the arXiv version.

\bibliographystyle{aasjournal}
\bibliography{methods}

\clearpage
\appendix
\section{The benchmark}
\label{app:bench}

Table~\ref{tab:bench} lists the \nmBenchN\ benchmark signals, their verdicts and
references, and the outcome of the ladder on the full HARPS-RVBank series of each host.

\begin{table*}
\caption{The benchmark: signals around HARPS-RVBank stars with later verdicts, and the
ladder on the full HARPS-RVBank series.}
\label{tab:bench}
\centering
\scriptsize
\setlength{\tabcolsep}{3pt}
\begin{tabular}{lrrcrrrrlcL{5.6cm}}
\hline\hline
Signal & $P$ (d) & $K$ ($\ms$) & Verdict & $N$ & $\Delta f\,T$ & $\log_{10}$FAP & $K/\sigma_K$ & Fails at & Trees & Discovery; verdict \\
\hline
\inputrows{tab_bench_rows}
\hline
\end{tabular}
\tablefoot{$P$, $K$: as claimed. Verdict: UI, upheld by independent evidence; UH, accepted
from HARPS data and not contested; R, refuted; D, disputed. $N$, $\log_{10}$FAP (global,
over 1.2--500\,d; 0.0 means a FAP near 1) and $K/\sigma_K$: this work, at the claimed
period. $\Delta f\,T$: offset of the evaluated frequency from the claimed one, in resolution
elements; $^e$ marks the edge of the $\pm0.5/T$ window. Fails at: the first rung of the a-priori rule the signal fails (det.: detection;
ind.: indicators; rot.: rotation; --: passes). Trees: score from the boosted trees, in
brackets where the signal is not significant and the score is outside the training domain.
References: discovery; later verdict. A single reference means the discovery paper itself
provides the verdict (e.g. transits) or, for UH, that no later paper contests it.}
\end{table*}


\section{Benchmark outcomes before and after the corrections}
\label{app:changes}

Table~\ref{tab:changes} lists the benchmark signals whose significance or first failed
rung changed between the frozen first-pass pipeline (rvvet 0.1) and this version.

\begin{table}
\caption{Benchmark signals whose outcome changed between the first pass (rvvet 0.1)
and this version.}
\label{tab:changes}
\centering
\footnotesize
\setlength{\tabcolsep}{3pt}
\begin{tabular}{lcrlrl}
\hline\hline
 & & \multicolumn{2}{c}{rvvet 0.1} & \multicolumn{2}{c}{rvvet \nmrvvetVersion} \\
Signal & Verdict & $\log_{10}$FAP & Fails at & $\log_{10}$FAP & Fails at \\
\hline
\inputrows{tab_v01_rows}
\hline
\end{tabular}
\tablefoot{Fails at: the first rung of the a-priori rule the signal fails (det.:
detection; ind.: indicators; rot.: rotation; --: passes).}
\end{table}
